%%% Statistical mechanics notes.

\documentclass{momento}

\usepackage{danielphysics}
\usepackage{danielfeynman}
\usepackage{tensor}

\title{Statistical Mechanics}
\author{Daniel Williams}


\providecommand{\Lag}{\mathcal{L}} %The Lagrangian
\providecommand{\Op}[1]{\mathrm{\hat{#1}}}  %A Quantum mechanical operator

\begin{document}
{
\thispagestyle{empty}
\begin{tikzpicture}[remember picture,overlay]
  \fill[fill=SpringGreen] (current page.south west) rectangle (current page.north east);
  %\fill[fill=white, yshift=-10cm]  (current page.north east) rectangle (current page.north west);
  \def\nbrcircles {377}
  \def\outerradius {30mm}
  \def\deviation {.9}
  \def\fudge {.62}

  \newcounter{cumulArea}
  \setcounter{cumulArea}{0}

  \pgfmathsetmacro {\goldenRatio} {(1+sqrt(5))}
  \pgfmathsetmacro {\meanArea} {pow(\outerradius * 10 / \nbrcircles, 2) * pi}
  \pgfmathsetmacro {\minArea} {\meanArea * (1 - \deviation)}
  \pgfmathsetmacro {\midArea} {\meanArea * (1 + \deviation) - \minArea}

  \foreach \b in {0,...,\nbrcircles}{
    % mod() must be used in order to calculate the right angle.
    % otherwise, when \b is greater than 28 the angle is greater
    % than 16384 and an error is raised ('Dimension too large').
    % -- thx Tonio for this one.
    \pgfmathsetmacro{\angle}{mod(\goldenRatio * \b, 2) * 180}

    \pgfmathsetmacro{\sratio}{\b / \nbrcircles}
    \pgfmathsetmacro{\smArea}{\minArea + \sratio * \midArea}
    \pgfmathsetmacro{\smRadius}{sqrt(\smArea / pi) / 2 * \fudge}
    \addtocounter{cumulArea}{\smArea};

    \pgfmathparse{sqrt(\value{cumulArea} / pi) / 2}
    \fill[opacity=0.3] (\angle:\pgfmathresult) circle [radius=\smRadius] ;
}  

  \node at (current page.center) [text width=11cm, yshift=10cm] 
    {\color{white}\fontsize{72pt}{80pt}\raggedright{ \selectfont\sffamily Statistical Mechanics}};

  \node at (current page.south) [text width= \textwidth, yshift=5cm] 
    {\color{white}\fontsize{32pt}{120pt}\selectfont \sffamily Daniel Williams};

\end{tikzpicture}
}
\newpage


\maketitle

\tableofcontents

\chapter{The Basic Postulates of Statistical Mechanics}
\label{cha:basic-post-stat}


\section{Definitions}
\label{sec:definitions}

\subsection{States}
\label{sec:states}

\begin{definition}[Macrostate]
  A macrostate is a description of a system which can be described by
  a few parameters, for example, temperature, pressure, internal
  energy, and volume.
\end{definition}

\begin{definition}[Microstate]
  A microstate is a detailed description of the constituent components of a macrostate, e.g.
  \begin{itemize}
  \item Classically, the position and momentum of each particle in the state;
  \item Quantum mechanically, the state vector for each particle.
  \end{itemize}
\end{definition}

\begin{postulate}
  \label{pos:averageover}
  We can identify the macroscopic parameters with the average at a
  given time over the set ensemble of microstates compatible with the
  macroscopic state and the environment of the system.
\end{postulate}
\begin{postulate}
  \label{pos:equalrep}
  For a given set of macroscopic parameters all quantum states
  consistent with this set have equal representation (i.e. weighting
  or probability) in the microcanonical ensemble.
\end{postulate}
\begin{postulate}[Averaging Postulate]
  All accessible microstates are equally probable.
\end{postulate}

\subsection{Ensembles}
\label{sec:ensembles}

\begin{definition}[Microcanonical Ensemble]
  The microcanonical ensemble represents the possible microstates of a
  mechanical system which have the same fixed total energy, and
  particle number.  All elements have the same volume, particle
  number, and total energy.
\end{definition}

\begin{definition}[Canonical Ensemble]
  This ensemble describes a closed system which is in contact with a
  heat bath. The total energy of each microstate is not constrained to
  be identical, with each system weighted by the Boltzmann factor. The
  volume, particle number, and temperature are fixed.
\end{definition}

\begin{definition}[Grand Canonical Ensemble]
  In this ensemble neither energy nor particle number are fixed, so
  temperature and chemical potential must be specified. This describes
  an open system.
\end{definition}

\subsection{Steps to model a system}
\label{sec:steps-model-system}

Problems in statistical modelling can be handled in broadly the same
way.
\begin{enumerate}
\item Solve the one-particle problem
\item Enumerate all of the possible distributions
\item Count the number of microstates which are in each distribution
\item Find the average distribution
\end{enumerate}

\subsection{Examples of Microcanonical Distributions}
\label{sec:exampl-micr-distr}

\begin{example}[Spin-1 Particle]
  Consider a single spin-1 particle in a magnetic field
  $\vec{B}$. There are three possible states, which have energies $E =
  \set{- \gamma m B}$ for $m=\set{-1, 0, 1}$.\\
  If $B=0$ then $E~{m}=0$ for all $m$, and so there is equal
  probability of being in any $m$ state, thus the probability of each
  state is $P = 1/3$. 
  If $\vec{B} \neq 0$ we cannot assign probabilities without a
  constraint;
  total energy, so the state is not microcanonical.\\
  If $E=0$ for $\vec{B}$ the system must have $m=0$ with $P=1$.
\end{example}
\begin{example}[Two spin-1 particles]
  Consider two spin-1 particles, $x$ and $y$ which are
  distinguishable. What is the probability of $m_x=+1$ while $m_y=-1$,
  i.e.~that the system has a state vector $(+1, -1)$?

  In the case that $B=0$ the state $(+1, -1)$ is one of 9 possible
  states, and so $P=\frac{1}{9}$. If, however, $B\neq 0$, we need to
  know $E$, so again, the state is non-microcanonical. If $E=0$ the
  states can be $(+1, -1)$, $(-1, +1)$, or $(0,0)$, so
  $P=\frac{1}{3}$.
\end{example}

\begin{example}[One dimensional polymer molecule]
  Consider a molecule consisting of $A$ links, each of length
  $b$. Each link may have one of two directions, denoted $\rightarrow$
  and $\leftarrow$. Each orientation has the same energy, so the
  ensemble is microcanonical, and this is an ideal microcanonical
  system.

  Let one end of the polymer be at $x=0$ and the other at $x = Lb$,
  for $-A \leq L \leq A$. Recalling postulate \ref{pos:averageover},
  the macroscopic length will be the average over all of the
  ensemble. From postulate \ref{pos:equalrep} we know that all of the
  $2^A$ microstates are equally probable, but this does not imply all
  macrostates are equally probable.

  A system of three links has $2^3$ possible configurations, so

  \begin{tabular}{cccc}
                                     & \rightarrow \leftarrow \leftarrow  & \leftarrow \rightarrow \rightarrow &                                     \\
    \leftarrow \leftarrow \leftarrow & \leftarrow \rightarrow \rightarrow & \rightarrow \leftarrow \rightarrow & \rightarrow \rightarrow \rightarrow \\ 
                                     & \leftarrow \leftarrow \rightarrow  & \rightarrow \rightarrow \leftarrow &                                     \\
$L = -3$                             & $L = -1$                           & $L =+1$                            & $L=+3$                              \\
$P=\frac{1}{8}$                      & $P=\frac{3}{8}$                    & $P=\frac{3}{8}$                    & $P=\frac{1}{8}$
  \end{tabular}

  Clearly there are a preferred set of orientations giving a length
  $|L|=1$ for the macroparameter. We can generalise this result by
  applying the binomial theorem.

  The total number of configurations with a fixed end at $x = Lb$ is
  given by a binomial distribution (see appendix
  \ref{sec:binom-distr}) such that
  \begin{equation}
    \label{eq:7}
    \Omega(L) = \frac{A!}{a_+! a_-!}
  \end{equation}
  where $a_+$ is the number of links pointing to the right, and $a_-$
  the number pointing to the left, such that 
  \[ a_+ + a_- = A, \quad a_+ - a_- = L \]
  Now,
  \begin{equation}
    \label{eq:8}
    \Omega(L) = \frac{A!}{\qty(\frac{A+L}{2})! \qty(\frac{A-L}{2})!}
  \end{equation}
  Then, taking Stirling's approximation (see appendix
  \ref{sec:stirl-appr}) 
  \begin{equation}
    \label{eq:9}
    P\qty(\frac{L}{A}) = \frac{\Omega(L)}{2^A} = \qty(\frac{2}{\pi a})^{\half} \exp( - \frac{L^2}{2A} )
  \end{equation}
  which has the form of a Gaussian distribution (which follows by the
  Central Limit Theorem). This indicates a most probable length of
  $L=0$.
\end{example}

\section{Canonical Distributions}
\label{sec:canon-distr}

% The microcanonical distribution affords great insight into physical
% systems but has limited direct relevance or application. We now
% introduce temperature, and remove the requirement that microstates all
% have the same energy. This allows the construction of a canonical
% ensemble from assemblies of microcanonical systems.  The system is
% controlled by contact with a heat bath at a temperature $T$, and all
% microstate energies, $E_i$ must be positive, so how are they
% distributed in probbaility?

In a microcanonical ensemble the energy of the state takes a
$\delta$-function form, but in a canonical ensemble, the energy of
each state is not the same---there is a distribution; we only know the
total energy of the system, which is constrained by the heat bath.

Although the microcanonical ensemble can be very useful it doesn't
occur often in real physical systems. A better approximation is
obtained by considering systems with a fixed number of particles,
volume, and temperature, held in a heat bath which defines the
temperature, $T$. The system is isolated, as the heat bath is
impermeable to particles, but energy is transferred to maintain the
temperature. (e.g. the average mark in a distribution of test results
being fixed; individuals can have a range of marks not equal to the
average, thus the systems within the canonical ensemble have an energy
constrained only by the average.)

We can build insight into a canonical system by building it from
smaller microcanonical systems which contribute overall to the
measurables. To see this, consider a system of $A$ identical
sub-systems sharing a total energy $E~{tot}$. Let $E_i$ denote the
energy of the $i$-th state. If $a_i$ is the number of systems at any
time $t$ with energy $E_i$ then the set of numbers $\set{a_i}$ satisfies 
\[ \sum_i a_i = A \]
and
\[ \sum_i a_i E_i = E~{tot} = AU = A \bar{E} \] for $\bar{E} = U$ the
average energy of the sub-systems.

Any set of $\set{a_i}$ satisfying these constraints represents a
possible mode of the distribution of total energy $E~{tot}$ among $A$
members of the ensemble. Any set $\set{a_i}$ satisfying the
constraints can be realised in a number of ways, e.g.~A reshuffle
among those members of the ensemble with different energy values, and
thus obtain a state of the ensemble which is distinct from the
original. How many ways are there to do this?

Let $\Omega$ be the number of ways that a set can be arranged, then
\begin{equation}
  \label{eq:1}
  \Omega(\set{a_i}) = \frac{A!}{a_1! a_2! a_3! \cdots} = \frac{A!}{\prod_i a_i!}
\end{equation}

Since all possible states of the ensemble are equally likely to occur
the frequency with which the distribution $\set{a_i}$ appears is
directly in proportion to $\Omega(\set{a_i})$. Thus, the most probable
mode of distribution is the one maximising $\Omega(\set{a_i})$, which
we denote $\set{a_i^{*}}$. This clearly satisfies the constraints, and
for all proactical purposes it's the only one which we need to
consider.

For large $A$ we expect $\Omega$ will be very strongly peaked, so
let's maximise $\Omega$, or, as it happens, maximise
$\frac{\log(\Omega)}{A}$, and define
\[ H = \frac{\log(\Omega)}{A} \]
We maximise $H$ subject to the constraints
\begin{subequations}
\begin{align}
  \sum a_i &= A \\
\sum a_i E_i &= E~{tot}
\end{align}
\end{subequations}

\begin{align*}
  H = \frac{\log(\Omega)}{A} &= \frac{1}{A} \log( \frac{A!}{a_1! a_2! \cdots}) \\
&= \frac{1}{A} \qty[ \log(A!) - \log(a_1! a_2! \cdots)]\\
&= \frac{1}{A} \qty[ A \log(A) - A - \floor{\sum_i a_i \log(a_i) - a_i}]
\end{align*}
Now we define the probability of being in state $a_i$ as 
\[ P_i = \frac{a_i}{A} \]
thus $\sum P_i = 1$.

So 
\begin{align*}
  A &= \frac{1}{A} \qty[ A \log(A) - A - \qty{ \sum_i A P_i \log(A P_i) - A P_i}] \\
&= \frac{1}{A} \qty[ A \log(A) - A - A \qty{ \sum_i P_i \qty[\log(A) + \log(P_i)] - P_i}]
\end{align*}
Cancellations mean that 
\[ H = - \sum P_i \log(P_i) \]
which needs to be maximised.

Let $\alpha$, $\beta$ be Lagrange multipliers, and 
\[ f = - \sum_i P_i \log(P_i) + \alpha(1 - \sum_i P_i) + \beta( u - \sum_i P_i E_i )\]
We then form the differential,
\[ \dd{f} = \sum_i \set{ - \log(P_i) - 1 - \alpha - \beta E_i}
\dd{P_i} = 0 \] This must hold for all values of $i$, so we can set
each side to equal $0$ independently,
\[ \therefore - \log(P_i) - 1 -\alpha - \beta E_i =0 \quad \forall i \]
\[ P_i = \exp( -1 -\alpha -\beta E_i) \]
and we also know $\sum P_i=1$, so
\begin{align*}
  \sum \exp(-1 -\alpha - \beta E_i) &= 1 \\
  e^{-(1+\alpha)} \sum e^{-\beta E_i} &= 1 \\
  e^{-(1+\alpha)} =  \qty(\sum e^{-\beta E_i})^{-1} &= \frac{1}{Z} 
\end{align*}
where $Z = \sum e^{-\beta E_i}$ is the partition function for the
system, the sum over all states weighted by the Boltzmann factor. Thus
\begin{equation}
  \label{eq:10}
  P_i = \frac{1}{Z} \exp(-\beta E_i)
\end{equation}
This can be generalised to reflect the fact that there are several
ways to reach the same energy state, and so we adopt the notation
\begin{equation}
  \label{eq:11}
  Z = \sum_i g_i \exp(- \beta E_i)
\end{equation}
for $g_i$ the multiplicity (or degeneracy) of the $i$th state.

The partition function, $Z$, is the central equation of statistical
mechanics, and knowledge of it allows the derivation of the major
results of thermodynamics.

\subsection{Major results using $Z$}
\label{sec:major-results-using}

\emph{The mean energy in a canonical ensemble} is given as
\[ \ev{E} = \sum P_i E_i = \sum \frac{1}{Z} E_i \exp(-\beta E_i) \]
Considering that 
\[ \pdv{\beta} \log(Z) = \frac{1}{Z} \pdv{\beta}(Z) \]
and 
\[ Z = Z = \sum_i \exp(- \beta E_i) \]
then
\[ \pdv{\beta} \log(Z) = - \frac{1}{Z} \sum_i E_i \exp(- \beta E_i) \]
and so
\begin{equation}
\label{eq:12}
 U = \ev{E} = - \pdv{\beta} \log(Z) 
\end{equation}

\emph{The energy fluctuations in a canonical ensemble} are
\begin{align*} 
\Delta E^2 &= \ev{E_i - \ev{E}}^2 = \ev{E_i^2 - 2 E_i \ev{E} + \ev{E}^2} \\
 &= \ev{E_i}^2 - \ev{E}^2 = \sum_i P_i E_i^2 - \qty( \sum_i P_i E_i )^2 \\
 &= \sum_i \frac{1}{Z} \exp(-\beta E_i) E_i^2 - \qty( \sum_i \frac{1}{Z} \exp(- \beta E_i) E_i )^2 \tag{\(\star\)}
\end{align*}
Noting that 
\begin{align*}
  - \pdv{U}{\beta} &= - \pdv{\beta}( \pdv{Z}{\beta} \frac{1}{Z} ) \\ 
&= \pdv{\beta} \qty[ \qty(\sum_i e^{-\beta E_i})^{-1} \sum_i \qty(-E_i e^{\beta E_i})] \\
&= \star
\end{align*}
then
\begin{equation}
  \label{eq:13}
  \Delta E^2 = - \pdv{U}{\beta}
\end{equation}
which is positive definite.

% \newpage

%%% Local Variables: 
%%% mode: latex
%%% TeX-master: "../project"
%%% End: 


\chapter{Relationships to Thermodynamics}
\label{cha:relat-therm}

\section{The First Law}
\label{sec:first-law}

\newcommand{\ddr}[1]{\dd{#1}}

We can reach a microscopic understanding of the First Law starting
with the Canonical Distribution,
\begin{equation}
  \label{eq:14}
  P_i = \frac{e^{-\beta E_i}}{Z}
\end{equation}
with the expression for $U = \ev{E} = \sum_i E_i P_i$. For a
reversible change of $E_i$ in a system due to a change in volume,
\begin{align*}
  \dd{u} &= \sum_i E_i \dd{P_i} + \sum_i P_i \dd{E_i} \\
&= \sum_i E_i \dd{P_i} + \sum_i P_i \qty( \pdv{E_i}{V} ) \dd{V} \\
&= \sum_i E_i \dd{P_i} - p \dd{V}
\end{align*}
This is equivalent to 
\begin{equation*}
  \dd{u} &= \ddr{Q} + \ddr{w}
\end{equation*}
for $\ddr{u} = \sum_i E_i \dd{P_i}$, and $\ddr{u}~{rev} = T \dd{S} =
\sum_iE_i \dd{P_i}$, for $S$ the entropy of a system.

\section{Relating $\beta$ and $T$}
\label{sec:relating-beta-t}

We can show that $\beta$ is a function of $T$ by onstruting a canonial
ensemble in which each system is a composite of two subsystems. These
must be in thermal equilibrium, and so the subsystems, $A$ and $B$,
must have the same temperature, $T_A = T_B$. We have the constraints 
\[ \sum_i a_{A,i} = A \qquad \sum_i a_{B,i} = A \]
Then
\[\sum a_{A,i} E_{A,i} + \sum_i a_{B,i} E_{B,i} = E~{tot} \]
The total energy in an ensemble is fixed but can be transferred
between subsystems. These three constraints imply three Lagrange
multipliers on the system, and we end up with
\[ P_{A,i} = \frac{e^{-\beta E_{A,i}}}{Z_A}, \qquad P_{B,i} =
\frac{e^{-\beta E_{B,i}}}{Z_B} \] where, as a result of the total
energy constraint $\beta$ is the same for each expression. If the two
subsystems always have the same $\beta$ and $T$ then 
\[ \beta = \beta(T) \]
Now, consider the function
\begin{equation}
  \label{eq:15}
  H = - \sum_i P_i \log(P_i)
\end{equation}
then, for a change,
\begin{align*}
  \dd{H} &= - \sum_i \dd{P_i} - \sum_i \dd{P_i} \log{P_i} \\
  &= - \sum_i \dd{P_i} -\sum_i \dd{P_i} \qty( log(e^{-\beta E_i}) - \log(z) ) \\
  &= - \sum_i \qty( \dd{P_i} \qty( - \beta E_i - \log(z) + \dd{P_i} )) \\
  &= - \sum_i \dd{P_i} \qty( 1 - \log(z) ) + \beta \sum \dd{P_i} E_i \\
  \text{since } & \sum P_i = 1 \implies \sum \dd{P_i} =0 \\
  &= \beta \sum_i \dd{P_i} E_i \\
  \shortintertext{then, assuming that $\beta T = \text{const} = k~B^{-1}$,}
  k~B \dd{H} &= \dd{s} \\
\therefore s &= k~B H = -k~B \sum_i P_i \log(P_i)
\end{align*}
This is an important but controversial result, but it is a consistent
form for $s$.

Consider a microcanonical system with a total of $\Omega$ states,
which are degenerate. Now assuming that the entropy is
\begin{equation}
  \label{eq:16}
  s = \phi(\Omega)
\end{equation}
and the probability of each microstate is proportional to $1/\Omega$,
taking two systems, $A$ and $B$ then the individual entropies are
\[ s~A = \phi(\Omega~A), \qquad s~B = \phi(\Omega~B) \] Entropy is
extensive, depending upon the total quantity of material in the
system, and so is additive for weakly interacting systems, so
\[ s~{AB} = s~A + s~B \] But the number of states must be
multiplicative, with $\Omega~{AB} = \Omega~A + \Omega~B$, so
\[ S~{AB} = \phi(\Omega~{AB}) = \phi(\Omega~A \Omega~B) =
\phi(\Omega~A) + \phi(\Omega~B) \] From this we can infer that the
relationship must be logarithmic, with
\begin{fequation}[Boltzmann entropy equation]
  \label{eq:17}
  s = k~B \log(\Omega)
\end{fequation}
Now we have $s$ as a function of $k~B H$:
\begin{align*}
  s &= -k~B \sum_i P_i \log(P_i) \\
&= - k~B \sum_i P_i \qty[ - \beta E_i - \log(Z) ]\\
&= k~B \beta \sum_i P_i E_i + k~B \log(z) \sum_i P_i
\intertext{since $\sum_i P_i =1$,}
&= \frac{1}{T} U + k~B \log(Z)
\end{align*}
Thus we have a new thermodynamic quantity,
\begin{fequation}[The Helmholtz free energy]
  \label{eq:18}
  F = U - Ts = -k~B T \log(Z)
\end{fequation}
which represents the amount of work which an isolated temperature at a
temperature $T$ can perform. This is the fundamental connection
between thermodynamics and statistical mechanics.

All of the other useful thermodynamic quantities can be derived using
this relationship, and its differential form,
\begin{equation}
  \label{eq:19}
  \dd{F} = - s \dd{T} - p \dd{V}.
\end{equation}

Thus
\begin{align}
  s = - \eval{\pdv{F}{T}}_V &= \qty( \pdv{T})_V \qty( k~B T \log(Z) ) \\
p = - \eval{\pdv{F}{V}}_T &= - k~B T \qty(\pdv{V})_T \qty( \log(Z) )
\end{align}
Recalling the Maxwell relations for $F=F(T,V,N)$, since $N$ can vary
in the grand canonical regime, we have 
\begin{equation}
  \label{eq:20}
  \dd{F} = - s \dd{T} - p \dd{V} + \mu \dd{N}
\end{equation}
for 
\begin{equation}
  \label{eq:21}
  \mu = \qty(\pdv{F}{N})_{T,V} = - k~B T \eval{ \pdv{V}}_{V,T} \log(Z)
\end{equation}
being the chemical potential.

Now consider a weakly interacting system of particles with a common
temperauture, $T$, then it is possible to define energy states of the
combined system in the form
\[ E~{Ai} + A~{Bj} \] which have a combined partition function,
\[ Z~{AB} = \sum_i \sum_j e^{- \beta E~{Ai} - \beta E~{Bj}} \] Now
consider $Z$ energies in each state, $A$ and $B$, for example,
\begin{align*}
  Z~{AB} &= \sum_{i=1}^2 \sum_{j=1}^2 e^{- \beta E~{Ai} - \beta E~{Bj}} \\
&= e^{-\beta A~1} e^{- \beta B_1} + e^{-\beta A_2} e^{-\beta B_1} \\ & \quad {}+ e^{-\beta A_1} e^{- \beta B_2} + e^{-\beta A_1}e^{-\beta B_2} \\
& = \qty( e^{-\beta E~A_1} + e^{-\beta E~{A_2}}) \qty( e^{-\beta E~{B_2}} + e^{-\beta E~{B_{2}}} ) \\ &= Z~A Z~B
\end{align*}
Thus the partition function of the combined weakly interacting system
is the product of the individual partition functions, and since
\begin{equation}
  \label{eq:22}
  F = - k~B T \log(Z) = -k~B T \log(Z~A Z~B) = F~A + F~B
\end{equation}
This multiplicative property of the partition functions is useful for
a range of complex systems, including the characterisation of
molecules' energies.

\section{Constructing a partition function}
\label{sec:constr-part-funct}

\subsection{A 2-energy-level system with $E_0=0$ and $E_1=E$	}
\label{sec:2-energy-level}

We have 
\[ Z = 1 + \exp(- \beta E) \] so the probability of each state being
occupied is
\begin{align*} P_0 &= \qty( 1+ \exp(-\beta E))^{-1}, \\ P_1 &= \exp(- \beta E) \qty( 1+ \exp(-\beta E))^{-1} \end{align*}
Thus
\[ F = -k~B T \log(1+ \exp(-\beta T)) \]

\subsection{Spin-$\half$ particle in a magnetic field}
\label{sec:spin-half-particle}

The energies of a spin-half particle in a magnetic field, $\vec{B}$,
will be \[ \pm E = \pm \gamma B \] corresponding to spin quantum
numbers $m \pm 1$, for $\gamma$ the projected magnetic moment. Thus
\begin{align*}
  Z &= \exp(\beta E) + \exp(- \beta E) \\ &= 2 \cosh(\beta E) \\
U &= F + Ts \\ s &= \pdv{T} (k~B T \log(Z) ) \\
U &= -k~B T \log(Z) - T \pdv{T} \qty( k~B T \log(Z)) \\ &= - k~B T \log(Z) + k~B T \log(Z) + k~B T^2 \pdv{T} \log(Z) \\ 
\pdv{T} \log(Z) &= \frac{1}{Z} \pdv{Z}{T}= \frac{1}{Z} \pdv{T} \qty(2 \cosh(\frac{E}{k~B T}) ) \\
&= \frac{2}{Z} \sinh( \frac{E}{k~B T}) \qty( - \frac{E}{k~B T^2} ) \\
&= - \frac{E}{k~B T^2} \tanh( \beta E) \\
U &= - E \tanh(\beta E) \\ &= - \gamma \beta \tanh(\frac{ \gamma \beta }{k~B T} )
\end{align*}
The mean magnetic moment can be found as
\begin{align*}
  \ev{m} &= \sum_i^2 P_i m_i = \gamma P_+ + (- \gamma) P_- \\
&= \gamma \qty[ \frac{1}{Z} \exp(\frac{\gamma B}{k~B T}) - \frac{1}{Z} \exp- \frac{\gamma B}{k~B T}] \\
&= 2 \frac{\gamma}{Z} \sinh( \frac{\gamma B}{k~B T} ) \\
&= \gamma \tanh( \frac{\gamma B}{k~B T})
\end{align*}
Thus
\begin{equation}
  \label{eq:1}
  U = - \ev{m} B
\end{equation}
which is what we expect. In the case that $\frac{\gamma B}{k~B T}$ is
very small, so in the case $T \gg 1$ or $B \ll 1$, then
\begin{fequation}[Curie Law]
  \label{eq:2}
  \ev{m} \approx \frac{\gamma^2 B}{k~B T} \sim \frac{1}{T}
\end{fequation}

\section{The harmonic oscillator}
\label{sec:harmonic-oscillator}

From both quantum mechanics and a simple analogy to standing waves in
a fixed volume, the energy of a fixed oscillator is 
\[ E_n = \qty(n+\half) h \nu, \quad \text{for}\ n \in \mathbb{N} \]
with $\nu$ the frequency of the oscillation. Then the partition
function can be found
\begin{align*}
  Z &= \sum^{\infty}_{n=0} \exp( - \frac{(n+\half) h \nu}{k~B T} )\\
&= \exp( - \half \frac{h \nu}{k~B T} ) \sum \exp(- \frac{h \nu}{k~B T})^n 
\intertext{Using the relation $\sum_{n=0}^{\infty} x^a = (1-x)^{-1}$}
&= \frac{\exp(- \half \frac{h \nu}{k~B T})}{1 - \exp( - \frac{h \nu}{k~B T}) }\\
&= \frac{\exp( - \frac{\Theta~\nu}{ 2 T})}{ 1- \exp(-\frac{\Theta~\nu}{T})}
\intertext{having defined the vibrational temperature,
\begin{equation}
  \label{eq:3}
  \Theta = \frac{h \nu}{k~B}
\end{equation}
}
Z &= \frac{1}{\exp( \frac{\Theta~\nu}{2 T}) - \exp( - \frac{\Theta_{\nu}}{2 T}) } \\
&= \frac{1}{2 \sinh( \frac{\Theta~\nu}{2 T} )}
\end{align*}
Armed with the partition function it is possible to find any
thermodynamic property of the system, for example
\[ U = k~B T^2 \dv{t} \log(Z) \]
So
\begin{align*}
  Z &= \exp(- \frac{\Theta~\nu}{2 T}) \qty( 1 - \exp(- \frac{\Theta~\nu}{T}) )^{-1} \\
\log(Z) &= - \frac{\Theta~\nu}{2 T} - \log( 1 - \exp(- \frac{\Theta~\nu}{T} ) ) \\
\dv{T} \log(Z) &= \frac{\Theta~\nu}{2T^2} - \frac{1}{1-\exp(\frac{\Theta~\nu}{T})} \qty[ - \qty( \frac{\Theta~\nu}{T}  ) \exp(- \frac{\Theta~\nu}{T}) ] \\
&= \frac{\Theta~\nu}{2 T^2} + \frac{\exp(- \frac{\Theta~\nu}{T})}{1 - \exp( \frac{\Theta~\nu}{T} )} \frac{\Theta~\nu}{T^2} \\
U &= k~B T^2 \dv{T} \log(Z) \\
&= \half k~B \Theta_{\nu} + k~B \Theta_\nu \frac{1}{\exp( \frac{\Theta~\nu}{T}) - 1}\\
&= \half h \nu + \frac{h \nu}{\exp( \frac{\Theta~\nu}{T}) -1}
\intertext{in the high temperature limit $\Theta~\nu/T \ll 1$}
U & \approx  \half h \nu + \frac{h \nu}{1 + \frac{h \nu}{k~B T} -1} \\
& \approx \half h \nu + k~B T \\
& \approx k~B T
\intertext{which is the classical two-dimensional result. \\In the low temperature limit we assume $\Theta \nu / T \gg 1$}
U & \approx \hal h \nu + h \nu \exp( - \frac{\Theta~\nu}{T}) 
\end{align*}
So as $\theta_{\nu}/T \to \infty$ we get
\begin{equation}
  \label{eq:4}
  U \approx \text{zero point energy}
\end{equation}
In classical thermodynamics the heat capacity, $C_{V}$, is constant,
that is $\dv{t} C_V = 0$ but here
\begin{align*}
  C_V &= \dv{U}{T} \approx h \nu \qty( \frac{\Theta_{\nu}}{T^2}) \exp(- \frac{\Theta_{\nu}}{T}) \\
&= \frac{h \nu \Theta_{\nu}}{T^2} \exp(- \frac{h \nu}{k~B T}) \to 0 \ \text{as} \ T \to 0
\end{align*}


%%% Local Variables: 
%%% mode: latex
%%% TeX-master: "../project"
%%% End: 


\part{Solids}
\label{part:solids}

\chapter{The Einstein Model}
\label{cha:einstein-model}

\part{Gases}
\label{part:gases}



\appendices

\chapter{Mathematical Techniques}
\label{cha:math-techn}


\section{Important Distributions}
\label{sec:import-distr}

\subsection{Binomial Distribution}
\label{sec:binom-distr}

For a system with two states the binomial distribution models the
number of one state compared to the other when taking $n$ samples with
replacement, given a total population $N$. So, for $n$ systems the
probability that $k$ are in one state is
\begin{equation}
  \label{eq:5}
  {n \choose k} = \frac{n!}{k!(n-k)!}
\end{equation}

\subsection{Multinomial Distribution}
\label{sec:mult-distr}

The multinomial distribution is a generalisation of the binomial
distribution in which a system may be in more than two possible
states.  The probability that $x_i$ systems are in state $i$, where
there are $p_i$ is the probability of a system being in state $i$, and
for $n$ systems, with $k$ possible states, then
\begin{equation}
  \label{eq:6}
  \begin{split}
    p(x_1, \dots, \x_k; n, p_1, \dots, p_k) = \Pr(X_1 = x_1, \dots, X_k = x_k) \\ = \frac{n!}{x_1!\cdots x_k!} p_1^{x_1}\cdots p_k^{x_k}
  \end{split}
\end{equation}

\section{Factorials}
\label{sec:factorials}

\subsection{Stirling's Approximation}
\label{sec:stirl-appr}

Stirling's Approximation is a means of approximating large factorials,
and is expressed as

\begin{equation}
  \label{eq:1}
  \log{n!} = n \log(n) - n + \mathcal{O}(log(n))
\end{equation}
which can be expressed more precisely as
\begin{equation}
  \label{eq:2}
  n! \approx \sqrt{2 \pi n} \qty( \frac{n}{e} )^n 
\end{equation}

\subsection{The gamma function}
\label{sec:gamma-function}

The gamma function is an extension of the factorial function to apply
to real and complex numbers. That is, if $n \in \mathbb{N}$,
\[\Gamma(n) = (n-1)! \]
If the real part of a complex number, $t$ is positive, then
\begin{equation}
  \label{eq:3}
  \Gamma(t) = \int_0^{\infty} x^{t-1} e^{-x} \dd{x}
\end{equation}
and has the factorial property,
\begin{equation}
  \label{eq:4}
  \Gamma(t+1) = t \Gamma(t)
\end{equation}

\section{The Method of Lagrange Multipliers}
\label{sec:meth-lagr-mult}

The method of Lagrange multipliers is used to optimise functions
subject to constraints. 

\begin{theorem}
  For any non-pathological function, $f(x_1, \dots, x_n)$, and for
  independent functions $g_1(x_1, \dots, x_n), \dots g_m(x_1, \dots,
  x_n)$, then, the local extrema of $f$, subject to these contraints
  can be found as the stationary points of the function
  \[ f + \lambda_1 g_1 + \lambda_2 g_2 + \cdots + \lambda_m g_m \]
  where $\lambda_1, \dots, \lambda_m$ are real scalars called Lagrange
  multipliers.
\end{theorem}

\begin{example}
  To find the extrema of 
\[ f(x,y) = x+y \]
on an ellipse \[ 4x^2 + y^2 = 80 \]

We then want the critical points of the function 
\begin{equation*} 
h(x, y) = f(x,y) + \lambda g(x,y) = x+y + \lambda (4 x^2 + y^2 -80) 
\end{equation*}
Thus $h_x= 1+8 \lambda x$ and $h_y = 1+ 2 \lambda y$, setting $h_x =
h_y =0$, then
\[ x = - \frac{1}{8} \lambda\ \quad y = - \half \lambda = 4x\] Thus
there is a critical point at $(x,y)=(x, 4x)$; substituting the
constraint,
\[ 4x^2 + (4x)^2 = 80 \implies x = \pm 2, y = \pm 8 \] So there are
critical points at $(2,8)$ and $(-2, -8)$, respectively a maximum and
a minimum.
\end{example}

Lagrange multipliers work on the basis that at a maximum, $f(x,y)$ of
$f$ we can not increase in any direction of a neighbouring point; if
we could we would move that way to a higher value, and not be at a
maximum. The extreme value of $f$ along the function of the constraint
can therefore only occur when the contours of the function $f$ are
parallel to the constraint. Gradients are perpendicular to the
contours, so if the function is parallel to the contours then the
gradients of both the constraint and the function must be either
parallel or antiparallel, so
\[ \nabla f = - \lambda \nabla g \] for a constant $\lambda$. Thus
$\nabla(f + \lambda g)=0$.


%%% Local Variables: 
%%% mode: latex
%%% TeX-master: "../project"
%%% End: 





\end{document}



%%% Local Variables: 
%%% mode: latex
%%% TeX-master: t
%%% End: 
